% ============================================================================
%  WSINDy on Manifold: Complete Mathematical Formulation and Code Reference
% ============================================================================
\documentclass[11pt]{article}

% ---- Packages ----
\usepackage[margin=1in]{geometry}
\usepackage{amsmath, amssymb, amsthm}
\usepackage{booktabs}
\usepackage{hyperref}
\usepackage{xcolor}
\usepackage{listings}
\usepackage{enumitem}
\usepackage[font=small]{caption}

% ---- Code listings ----
\definecolor{codebg}{gray}{0.96}
\definecolor{codegreen}{rgb}{0.1,0.5,0.1}
\definecolor{codeblue}{rgb}{0.1,0.1,0.6}
\definecolor{codegray}{rgb}{0.4,0.4,0.4}
\lstset{
  language=Python,
  basicstyle=\ttfamily\small,
  backgroundcolor=\color{codebg},
  keywordstyle=\color{codeblue}\bfseries,
  commentstyle=\color{codegreen}\itshape,
  stringstyle=\color{codegray},
  numbers=left,
  numberstyle=\tiny\color{codegray},
  numbersep=5pt,
  frame=single,
  framerule=0.4pt,
  breaklines=true,
  breakatwhitespace=false,
  showstringspaces=false,
  tabsize=4,
  xleftmargin=2em,
  framexleftmargin=1.5em,
  captionpos=b,
  escapeinside={(*@}{@*)},
}

% ---- Custom commands ----
\newcommand{\R}{\mathbb{R}}
\newcommand{\bx}{\mathbf{x}}
\newcommand{\bp}{\mathbf{p}}
\newcommand{\bv}{\mathbf{v}}
\DeclareMathOperator{\Div}{div}
\DeclareMathOperator{\diag}{diag}
\DeclareMathOperator*{\argmin}{arg\,min}
\newcommand{\colvec}[1]{\begin{pmatrix}#1\end{pmatrix}}

\title{%
  \textbf{Weak SINDy on Manifold}\\[6pt]
  \large Complete Mathematical Formulation and Code Reference\\
  \large for Multi-Field PDE Discovery of Vicsek--Morse Crowd Dynamics}
\author{Macias-Olivares, E.}
\date{February 2026}

\begin{document}
\maketitle
\tableofcontents
\newpage

% ============================================================================
\section{Overview and Module Map}
% ============================================================================

This document provides the complete mathematical derivation and code reference
for the WSINDy (Weak Sparse Identification of Nonlinear Dynamics) package
in \texttt{src/wsindy/}. The codebase discovers governing PDEs from
spatiotemporal data on 2D periodic domains by:

\begin{enumerate}[leftmargin=*]
\item \textbf{Computing fields} (density $\rho$, flux $\bp$, Morse potential $\Phi$) from particle data;
\item \textbf{Constructing a weak-form linear system} $\mathbf{b} = \mathbf{G}\mathbf{w}$
  by convolving candidate library terms with compactly-supported test functions;
\item \textbf{Fitting sparse models} via Modified Sequential Threshold Least Squares (MSTLS);
\item \textbf{Forecasting} the discovered PDE using RK4 or ETDRK4 time integrators.
\end{enumerate}

\subsection{Module Structure}

\begin{center}
\begin{tabular}{ll}
\toprule
\textbf{Module} & \textbf{Purpose} \\
\midrule
\texttt{grid.py}            & \texttt{GridSpec} dataclass ($\Delta t, \Delta x, \Delta y$) \\
\texttt{fields.py}          & Field computation: KDE flux, Morse potential, derivatives \\
\texttt{operators.py}       & Spectral derivative utilities (standalone) \\
\texttt{test\_functions.py} & Separable test function $\psi$ and its derivatives \\
\texttt{system.py}          & Weak system builder: query points, FFT convolution \\
\texttt{library.py}         & \texttt{LibraryBuilder}: polynomial and nonlinear terms \\
\texttt{multifield.py}      & 3-equation system: library terms, discovery, forecast \\
\texttt{regression.py}      & MSTLS: preconditioning, dominant-balance thresholding \\
\texttt{fit.py}             & High-level fitting: precondition $\to$ MSTLS $\to$ unscale \\
\texttt{select.py}          & Model selection over test-function scales $\ell$ \\
\texttt{model.py}           & \texttt{WSINDyModel} dataclass \\
\texttt{integrators.py}     & RK4 and ETDRK4 time integrators \\
\texttt{forecast.py}        & High-level forecast dispatch \\
\texttt{rhs.py}             & Strong-form RHS evaluation \\
\texttt{uncertainty.py}     & Bootstrap UQ for coefficient distributions \\
\bottomrule
\end{tabular}
\end{center}


% ============================================================================
\section{Grid Specification}
\label{sec:grid}
% ============================================================================

The \texttt{GridSpec} dataclass (in \texttt{grid.py}) stores the fundamental discretisation:

\begin{lstlisting}[caption={grid.py --- GridSpec dataclass}]
@dataclass(frozen=True)
class GridSpec:
    dt: float
    dx: float
    dy: float
    periodic_space: bool = True
    periodic_time: bool = False
\end{lstlisting}

All spatiotemporal data lives on a uniform grid
$\Omega = [0, L_x) \times [0, L_y)$ with $N_x \times N_y$ cells and
temporal spacing $\Delta t$.  Both spatial directions are periodic;
time is not.


% ============================================================================
\section{Field Computation}
\label{sec:fields}
% ============================================================================

The three primary fields are the density $\rho(x,y,t)$, the
polarisation/flux $\bp(x,y,t) = (p_x, p_y)$, and the Morse potential
$\Phi(x,y,t)$.  All are defined on the same $(T, N_y, N_x)$ array.


\subsection{Density via Histogram + Gaussian Smoothing}

Particle positions $\{\bx_i(t)\}_{i=1}^N$ are binned into a 2D
histogram $H^t \in \R^{N_y \times N_x}$ and smoothed:
\begin{equation}
  \rho(x,y,t) \;=\; \bigl(G_\sigma * H^t\bigr)(x,y),
  \qquad
  H^t_{jl} = \frac{1}{\Delta x \,\Delta y}
  \sum_{i=1}^N \mathbf{1}[\bx_i(t) \in C_{jl}],
\end{equation}
where $G_\sigma$ is a 2D Gaussian kernel with bandwidth $\sigma$ (in
grid-cell units) and the convolution uses periodic wrapping.


\subsection{Flux / Polarisation via Velocity-Weighted KDE}

The polarisation field encodes both density and local average velocity:
\begin{equation}
  \bp(x,t) = \sum_{i=1}^N K_\sigma\!\bigl(x - \bx_i(t)\bigr)\,\bv_i(t),
\end{equation}
i.e., the same Gaussian kernel as the density KDE, but weighted by each
particle's velocity $\bv_i$.  In practice we compute this per-component:

\begin{lstlisting}[caption={fields.py --- compute\_flux\_kde (key loop)}]
for out_idx, t in enumerate(range(0, T_full, subsample)):
    pos_t = traj[t]   # (N, 2)
    vel_t = vel[t]     # (N, 2)

    # Weighted histograms
    hx, _, _ = np.histogram2d(
        pos_t[:, 0], pos_t[:, 1],
        bins=[x_edges, y_edges],
        weights=vel_t[:, 0],
    )
    hy, _, _ = np.histogram2d(
        pos_t[:, 0], pos_t[:, 1],
        bins=[x_edges, y_edges],
        weights=vel_t[:, 1],
    )

    # Convert to flux density (per unit area)
    hx = hx.T / (dx_phys * dy_phys)
    hy = hy.T / (dx_phys * dy_phys)

    # Gaussian smoothing (same kernel as density)
    if bandwidth > 0:
        hx = gaussian_filter(hx, sigma=bandwidth, mode=mode)
        hy = gaussian_filter(hy, sigma=bandwidth, mode=mode)

    px[out_idx] = hx
    py[out_idx] = hy
\end{lstlisting}


\subsection{Morse Potential via FFT Convolution}

The Morse interaction potential creates short-range repulsion and
long-range attraction:
\begin{equation}
  W(r) = -C_r \exp\!\bigl(-r/l_r\bigr) + C_a \exp\!\bigl(-r/l_a\bigr),
  \quad
  \Phi(x,t) = (W * \rho)(x,t) = \int_\Omega W(|x - x'|)\,\rho(x',t)\,dx'.
\end{equation}

The self-interaction is removed: $W(0) = 0$.
On the periodic domain, convolution is computed efficiently via the
convolution theorem:
\begin{equation}
  \Phi(x,t) = \mathcal{F}^{-1}\!\bigl[\hat{W}\cdot\hat{\rho}(\cdot,t)\bigr](x)
  \;\cdot\;\Delta x\,\Delta y.
\end{equation}

\begin{lstlisting}[caption={fields.py --- Morse kernel construction}]
def morse_kernel_grid(xgrid, ygrid, Lx, Ly, Cr, Ca, lr, la, bc="periodic"):
    nx, ny = len(xgrid), len(ygrid)
    # FFT-centered distances (periodic wrapping)
    cx = np.fft.fftfreq(nx, d=1.0) * Lx
    cy = np.fft.fftfreq(ny, d=1.0) * Ly
    CX, CY = np.meshgrid(cx, cy)
    R = np.sqrt(CX**2 + CY**2)
    R = np.maximum(R, 1e-10)
    W = -Cr * np.exp(-R / lr) + Ca * np.exp(-R / la)
    W[0, 0] = 0.0   # remove self-interaction
    return W
\end{lstlisting}

\begin{lstlisting}[caption={fields.py --- FFT convolution for $\Phi = W * \rho$}]
def compute_morse_potential(rho, xgrid, ygrid, Lx, Ly, Cr, Ca, lr, la, bc="periodic"):
    W = morse_kernel_grid(xgrid, ygrid, Lx, Ly, Cr, Ca, lr, la, bc)
    dx, dy = Lx / len(xgrid), Ly / len(ygrid)
    W_hat = np.fft.fft2(W)
    Phi = np.zeros_like(rho)
    for t in range(rho.shape[0]):
        rho_hat = np.fft.fft2(rho[t])
        Phi[t] = np.real(np.fft.ifft2(W_hat * rho_hat)) * dx * dy
    return Phi
\end{lstlisting}


\subsection{Spatial Derivatives: FD and Spectral Backends}

The \texttt{FieldData} class computes all spatial derivatives via a
switchable backend.  By default, training uses finite differences (FD);
forecasting uses spectral (FFT) derivatives.

\paragraph{Finite-Difference Backend.}
Central differences with periodic wrapping:
\begin{equation}
  \frac{\partial f}{\partial x}\bigg|_{j}
  \approx \frac{f_{j+1} - f_{j-1}}{2\,\Delta x},
  \qquad
  \frac{\partial^2 f}{\partial x^2}\bigg|_{j}
  \approx \frac{f_{j+1} - 2f_j + f_{j-1}}{\Delta x^2}.
\end{equation}

The Laplacian is $\Delta f = \partial_{xx} f + \partial_{yy} f$ and
the biharmonic is $\Delta^2 f = \Delta(\Delta f)$.

\begin{lstlisting}[caption={fields.py --- Finite-difference stencils}]
def _dx(f, dx):
    return (np.roll(f, -1, axis=-1) - np.roll(f, 1, axis=-1)) / (2 * dx)

def _dy(f, dy):
    return (np.roll(f, -1, axis=-2) - np.roll(f, 1, axis=-2)) / (2 * dy)

def _dxx(f, dx):
    return (np.roll(f, -1, axis=-1) - 2*f + np.roll(f, 1, axis=-1)) / dx**2

def _dyy(f, dy):
    return (np.roll(f, -1, axis=-2) - 2*f + np.roll(f, 1, axis=-2)) / dy**2

def _lap(f, dx, dy):
    return _dxx(f, dx) + _dyy(f, dy)

def _bilap(f, dx, dy):
    return _lap(_lap(f, dx, dy), dx, dy)
\end{lstlisting}

\paragraph{Spectral (FFT) Backend.}
Machine-precision on periodic grids, unconditionally stable:
\begin{equation}
  \frac{\partial f}{\partial x}
  = \mathcal{F}^{-1}\bigl[ik_x\,\hat{f}\bigr],
  \qquad
  \Delta f = \mathcal{F}^{-1}\bigl[-(k_x^2 + k_y^2)\,\hat{f}\bigr],
\end{equation}
where $k_x = 2\pi \cdot \text{fftfreq}(N_x, \Delta x)$ and similarly
for $k_y$.  The Nyquist mode is zeroed for odd-order derivatives.

\begin{lstlisting}[caption={fields.py --- Spectral derivatives}]
def _spectral_dx(f, Lx):
    nx = f.shape[-1]
    dx = Lx / nx
    kx = 2 * np.pi * np.fft.fftfreq(nx, d=dx)
    if nx % 2 == 0:
        kx[nx // 2] = 0.0   # zero Nyquist
    f_hat = np.fft.fft(f, axis=-1)
    return np.real(np.fft.ifft(1j * kx * f_hat, axis=-1))

def _spectral_lap(f, Lx, Ly):
    ny, nx = f.shape[-2], f.shape[-1]
    K2 = _spectral_wavenumbers_2d(ny, nx, Lx, Ly)
    f_hat = np.fft.fft2(f, axes=(-2, -1))
    return np.real(np.fft.ifft2(-K2 * f_hat, axes=(-2, -1)))
\end{lstlisting}


\subsection{The FieldData Container}

All derived quantities (divergences, Laplacians, nonlinear terms) are
lazily computed and cached inside \texttt{FieldData}:

\begin{lstlisting}[caption={fields.py --- FieldData with caching dispatch}]
class FieldData:
    def __init__(self, rho, px, py, grid, Lx, Ly,
                 Phi=None, use_spectral=False):
        self.rho = rho; self.px = px; self.py = py
        self.Phi = Phi; self.grid = grid
        self.Lx = Lx; self.Ly = Ly
        self.use_spectral = use_spectral

    def _ddx(self, f):
        if self.use_spectral:
            return _spectral_dx(f, self.Lx)
        return _dx(f, self.dx)

    def _laplacian(self, f):
        if self.use_spectral:
            return _spectral_lap(f, self.Lx, self.Ly)
        return _lap(f, self.dx, self.dy)

    # Cached derived fields
    def div_p(self):
        return self._cache("div_p",
            lambda: self._ddx(self.px) + self._ddy(self.py))

    def rho_grad_Phi_x(self):
        return self._cache("rho_grad_Phi_x",
            lambda: self.rho * self._ddx(self.Phi))

    def div_rho_gradPhi(self):
        return self._cache("div_rho_gradPhi", lambda: self._div(
            self.rho_grad_Phi_x(), self.rho_grad_Phi_y()))

    def p_dot_grad_px(self):
        return self._cache("p_dot_grad_px", lambda:
            self.px * self._ddx(self.px) + self.py * self._ddy(self.px))
\end{lstlisting}

The full list of cached derived fields:
\begin{itemize}[leftmargin=*,itemsep=1pt]
  \item \textbf{Density:} $\partial_x\rho$, $\partial_y\rho$, $\Delta\rho$, $\rho^2$, $\rho^3$
  \item \textbf{Flux:} $\nabla\cdot\bp$, $\Delta p_x$, $\Delta p_y$, $\Delta^2 p_x$, $\Delta^2 p_y$, $|\bp|^2$
  \item \textbf{Nonlinear density:} $\Delta(\rho^2)$, $\Delta(\rho^3)$, $\Delta(|\bp|^2)$
  \item \textbf{Morse:} $\nabla\Phi$, $\rho\nabla\Phi$, $\nabla\cdot(\rho\nabla\Phi)$
  \item \textbf{Self-advection:} $(\bp\cdot\nabla)p_x$, $(\bp\cdot\nabla)p_y$
  \item \textbf{Pressure gradients:} $\partial_x(\rho^2)$, $\partial_y(\rho^2)$
\end{itemize}


% ============================================================================
\section{Test Functions}
\label{sec:test_functions}
% ============================================================================

The weak formulation transfers derivatives from the data onto smooth,
compactly supported test functions $\psi$, dramatically improving noise
robustness.

\subsection{Separable Polynomial Bump}

We build $\psi$ as a product of 1-D polynomial bumps:
\begin{equation}\label{eq:psi_separable}
  \psi(x,y,t) = \varphi_t(t)\;\varphi_x(x)\;\varphi_y(y),
\end{equation}
where each 1-D factor is a compactly supported polynomial bump:
\begin{equation}\label{eq:bump}
  \varphi_i(s)
  = \biggl(1 - \frac{s^2}{\ell_i^2\,\Delta_i^2}\biggr)^{p_i}_{\!+},
  \qquad |s| \le \ell_i\,\Delta_i,
\end{equation}
supported on $2\ell_i + 1$ grid points in direction $i$.  The
parameters $\ell = (\ell_t, \ell_x, \ell_y)$ control the test-function
width (number of half-width grid points) and
$p = (p_t, p_x, p_y)$ control the polynomial exponent (smoothness).

\begin{lstlisting}[caption={test\_functions.py --- 1-D bump construction}]
def _bump_1d(n_half, p, ds):
    """Polynomial bump: (1 - (s/(n_half*ds))^2)^p on [-n_half..n_half]."""
    s = np.arange(-n_half, n_half + 1, dtype=np.float64) * ds
    width = n_half * ds
    phi = np.maximum(1.0 - (s / width)**2, 0.0)**p
    phi /= (phi.sum() * ds)  # normalise
    return phi
\end{lstlisting}


\subsection{Derivative Bundle}

The weak system requires the following test-function derivatives,
all computed analytically from the product rule:
\begin{align}
  \psi_t &= \varphi_t'(t)\;\varphi_x(x)\;\varphi_y(y), \\
  \psi_x &= \varphi_t(t)\;\varphi_x'(x)\;\varphi_y(y), \\
  \psi_y &= \varphi_t(t)\;\varphi_x(x)\;\varphi_y'(y), \\
  \psi_{xx} &= \varphi_t(t)\;\varphi_x''(x)\;\varphi_y(y), \\
  \psi_{yy} &= \varphi_t(t)\;\varphi_x(x)\;\varphi_y''(y).
\end{align}

The 1-D derivatives are computed by central finite differences on the
bump's own support lattice.

\begin{lstlisting}[caption={test\_functions.py --- make\_separable\_psi}]
def make_separable_psi(grid, ellt, ellx, elly, pt=2, px=2, py=2):
    """Build the test-function bundle psi and all derivatives."""
    phi_t = _bump_1d(ellt, pt, grid.dt)
    phi_x = _bump_1d(ellx, px, grid.dx)
    phi_y = _bump_1d(elly, py, grid.dy)

    # 3-D separable product via outer products
    psi = np.einsum('i,j,k->ijk', phi_t, phi_x, phi_y)

    # Derivatives via product rule + central FD on the 1-D bumps
    dphi_t = np.gradient(phi_t, grid.dt)
    dphi_x = np.gradient(phi_x, grid.dx)
    dphi_y = np.gradient(phi_y, grid.dy)
    d2phi_x = np.gradient(dphi_x, grid.dx)
    d2phi_y = np.gradient(dphi_y, grid.dy)

    psi_t  = np.einsum('i,j,k->ijk', dphi_t,  phi_x,   phi_y)
    psi_x  = np.einsum('i,j,k->ijk', phi_t,   dphi_x,  phi_y)
    psi_y  = np.einsum('i,j,k->ijk', phi_t,   phi_x,   dphi_y)
    psi_xx = np.einsum('i,j,k->ijk', phi_t,   d2phi_x, phi_y)
    psi_yy = np.einsum('i,j,k->ijk', phi_t,   phi_x,   d2phi_y)

    return {"psi": psi, "psi_t": psi_t, "psi_x": psi_x,
            "psi_y": psi_y, "psi_xx": psi_xx, "psi_yy": psi_yy,
            "phi_t": phi_t, "phi_x": phi_x, "phi_y": phi_y}
\end{lstlisting}


% ============================================================================
\section{Weak Formulation}
\label{sec:weak}
% ============================================================================

\subsection{From Strong to Weak Form}

Consider a PDE of the form:
\begin{equation}\label{eq:strong}
  u_t = \sum_{m=1}^M w_m\, \mathcal{D}_m[f_m(u)],
\end{equation}
where $\mathcal{D}_m$ is a spatial differential operator (identity,
$\partial_x$, $\partial_y$, $\Delta$, etc.) and $f_m$ is a nonlinear
feature ($u, u^2, u^3, \ldots$).

Multiplying both sides by a test function $\psi$ and integrating by parts
(exploiting compact support of $\psi$ and periodic boundary conditions):
\begin{equation}\label{eq:weak_form}
  \underbrace{\int_\Omega \psi_t \cdot u \;dx\,dt}_{= b}
  \;=\; \sum_{m=1}^M w_m
  \underbrace{\int_\Omega \psi^{(m)} \cdot f_m(u) \;dx\,dt}_{= G_m},
\end{equation}
where $\psi^{(m)}$ is the \emph{adjoint kernel}: integration by parts
moves derivatives from $u$ onto $\psi$.  The kernel map is:
\begin{center}
\begin{tabular}{ll}
\toprule
Operator $\mathcal{D}_m$ & Kernel $\psi^{(m)}$ used \\
\midrule
$I$ (identity) & $\psi$ \\
$\partial_x$   & $-\psi_x$ \\
$\partial_y$   & $-\psi_y$ \\
$\Delta$ (Laplacian) & $\psi_{xx} + \psi_{yy}$ \\
\bottomrule
\end{tabular}
\end{center}

\begin{lstlisting}[caption={system.py --- Kernel mapping}]
def get_kernel(op, psi_bundle):
    """Map operator name to the appropriate psi derivative."""
    if op == "I":
        return psi_bundle["psi"]
    elif op == "dx":
        return -psi_bundle["psi_x"]
    elif op == "dy":
        return -psi_bundle["psi_y"]
    elif op == "lap":
        return psi_bundle["psi_xx"] + psi_bundle["psi_yy"]
    raise ValueError(f"Unknown operator: {op}")
\end{lstlisting}


\subsection{Query Points}

The integrals are discretised by sampling at \emph{query points}
$(t_q, x_q, y_q)$  on a uniform stride:
\begin{equation}
  b_q = (\psi_t * u)(t_q, x_q, y_q),
  \qquad
  G_{q,m} = (\psi^{(m)} * f_m(u))(t_q, x_q, y_q).
\end{equation}

Query points avoid the temporal boundaries (a margin of $\ell_t + 1$
points on each side) to ensure the test function support lies within
the data.

\begin{lstlisting}[caption={system.py --- Query point generation}]
def make_query_indices(T, nx, ny, stride_t=2, stride_x=2, stride_y=2,
                       t_margin=None):
    """Uniform strided (t,x,y) query points with time margins."""
    if t_margin is None:
        t_margin = 3
    t_idx = np.arange(t_margin, T - t_margin, stride_t)
    x_idx = np.arange(0, nx, stride_x)
    y_idx = np.arange(0, ny, stride_y)
    # Build (K, 3) array of all (t, x, y) tuples
    tt, xx, yy = np.meshgrid(t_idx, x_idx, y_idx, indexing='ij')
    return np.column_stack([tt.ravel(), xx.ravel(), yy.ravel()])
\end{lstlisting}


\subsection{FFT Convolution}

The convolutions $(\psi * u)$ are computed via 3-D FFT, with periodic
boundary conditions in space and non-periodic (zero-padded) in time:

\begin{lstlisting}[caption={system.py --- FFT convolution (3D)}]
def fft_convolve3d_same(data, kernel, periodic=(False, True, True)):
    """3-D FFT convolution returning same-size output.

    Non-periodic axes are zero-padded to avoid wrap-around.
    Periodic axes use circular convolution (direct FFT).
    """
    shape_out = data.shape
    pad_shape = []
    for i in range(3):
        if periodic[i]:
            pad_shape.append(data.shape[i])  # no padding needed
        else:
            pad_shape.append(data.shape[i] + kernel.shape[i] - 1)

    D = np.fft.fftn(data, s=pad_shape)
    K = np.fft.fftn(kernel, s=pad_shape)
    conv = np.real(np.fft.ifftn(D * K))

    # Crop to original size (centered for non-periodic axes)
    slices = []
    for i in range(3):
        if periodic[i]:
            slices.append(slice(0, shape_out[i]))
        else:
            offset = kernel.shape[i] // 2
            slices.append(slice(offset, offset + shape_out[i]))
    return conv[tuple(slices)]
\end{lstlisting}


\subsection{Building the Weak System}

The full system $\mathbf{b} = \mathbf{G}\mathbf{w}$ is assembled as:

\begin{lstlisting}[caption={system.py --- build\_weak\_system}]
def build_weak_system(U, grid, psi_bundle, library_terms, query_idx):
    """Assemble b = G w for the scalar weak-form system.

    U : (T, nx, ny) spatiotemporal data
    library_terms : list of (op, feature) tuples
    query_idx : (K, 3) array
    """
    qt, qx, qy = query_idx[:, 0], query_idx[:, 1], query_idx[:, 2]
    K, M = len(qt), len(library_terms)
    periodic = (False, True, True)

    # LHS: b = conv(psi_t, U) at query points
    psi_t = psi_bundle["psi_t"]
    b_full = fft_convolve3d_same(U, psi_t, periodic)
    b = b_full[qt, qx, qy]

    # RHS: G[:,m] = conv(kernel_m, f_m(U)) at query points
    G = np.empty((K, M))
    col_names = []
    for m, (op, feat) in enumerate(library_terms):
        kernel = get_kernel(op, psi_bundle)
        f_U = eval_feature(U, feat)  # e.g., U, U**2, U**3
        conv_m = fft_convolve3d_same(f_U, kernel, periodic)
        G[:, m] = conv_m[qt, qx, qy]
        col_names.append(f"{op}:{feat}")

    return b, G, col_names
\end{lstlisting}


% ============================================================================
\section{Candidate Library}
\label{sec:library}
% ============================================================================

\subsection{Scalar Library Builder}

The \texttt{LibraryBuilder} (in \texttt{library.py}) constructs the set of
candidate terms as $(operator, feature)$ pairs:

\begin{lstlisting}[caption={library.py --- LibraryBuilder usage}]
builder = LibraryBuilder()
builder.add_polynomials(max_power=3, operators=["I","dx","dy","lap"])
# Generates: I:u, I:u2, I:u3, dx:u, dx:u2, ..., lap:u, lap:u2, lap:u3
# Excludes nonsensical terms like dx:1

builder.add_nonlinear_diffusion(powers=[2, 3])
# Adds: lap:u2, lap:u3  (if not already present)

builder.add_gradient_nonlinearity()
# Adds: I:|nabla u|^2
\end{lstlisting}

These map to the mathematical library:
\begin{align}
  &\text{Polynomials:} &&
  u,\; u^2,\; u^3,\;
  \partial_x u,\; \partial_y u,\;
  \partial_x(u^2),\; \ldots,\;
  \Delta u,\; \Delta(u^2),\; \Delta(u^3). \\
  &\text{Gradient nonlinearity:} &&
  |\nabla u|^2 = (\partial_x u)^2 + (\partial_y u)^2.
\end{align}


\subsection{Multi-Field Library (3-Equation System)}
\label{sec:multifield_library}

For the coupled Vicsek--Morse system, each equation gets its own
library of \texttt{LibraryTerm} objects.  Each term directly provides
a $(T, N_y, N_x)$ field from \texttt{FieldData}.

\paragraph{$\rho$-equation library:}
\begin{equation}
  \rho_t = c_1\,\nabla\cdot\bp
  + c_2\,\Delta\rho
  + c_3\,\nabla\cdot(\rho\nabla\Phi)
  + c_4\,\Delta(\rho^2)
  + c_5\,\Delta(\rho^3)
  + \ldots
\end{equation}

\begin{center}
\begin{tabular}{lll}
\toprule
\textbf{Code Name} & \textbf{Math} & \textbf{Category} \\
\midrule
\texttt{div\_p} & $\nabla\cdot\bp$ & C1 (continuity) \\
\texttt{lap\_rho} & $\Delta\rho$ & C1 (diffusion) \\
\texttt{div\_rho\_gradPhi} & $\nabla\cdot(\rho\nabla\Phi)$ & C2 (Morse aggregation) \\
\texttt{lap\_rho2} & $\Delta(\rho^2)$ & C3 (nonlinear diffusion) \\
\texttt{lap\_rho3} & $\Delta(\rho^3)$ & C3 (crowd pressure) \\
\texttt{lap\_p\_sq} & $\Delta(|\bp|^2)$ & C4 (coupling, dangerous) \\
\bottomrule
\end{tabular}
\end{center}

\paragraph{$p_x$-equation library:}
\begin{equation}
  (p_x)_t = d_1\, p_x
  + d_2\,|\bp|^2 p_x
  + d_3\,\partial_x\rho
  + d_4\,\Delta p_x
  + d_5\,\rho\,\partial_x\Phi
  + \ldots
\end{equation}

\begin{center}
\begin{tabular}{lll}
\toprule
\textbf{Code Name} & \textbf{Math} & \textbf{Category} \\
\midrule
\texttt{px} & $p_x$ & D1 (linear alignment) \\
\texttt{p\_sq\_px} & $|\bp|^2\,p_x$ & D1 (saturation) \\
\texttt{dx\_rho} & $\partial_x\rho$ & D2 (pressure) \\
\texttt{lap\_px} & $\Delta p_x$ & D3 (viscosity) \\
\texttt{rho\_dx\_Phi} & $\rho\,\partial_x\Phi$ & D5 (Morse forcing) \\
\texttt{p\_dot\_grad\_px} & $(\bp\cdot\nabla)p_x$ & D4 (self-advection) \\
\texttt{bilap\_px} & $\Delta^2 p_x$ & D3+ (hyper-viscosity) \\
\texttt{dx\_rho2} & $\partial_x(\rho^2)$ & D2+ (extended pressure) \\
\bottomrule
\end{tabular}
\end{center}

The $p_y$ library is symmetric, replacing $\partial_x \to \partial_y$.

\begin{lstlisting}[caption={multifield.py --- Library term example}]
class LibraryTerm:
    def __init__(self, name, evaluator, *, latex="", equation="rho",
                 dangerous=False):
        self.name = name
        self.evaluator = evaluator  # f(FieldData) -> (T, ny, nx)
        self.latex = latex
        self.equation = equation
        self.dangerous = dangerous

# Example term factories
def rho_terms_core():
    return [
        LibraryTerm("div_p", lambda fd: fd.div_p(),
                    latex=r"\nabla\cdot\mathbf{p}", equation="rho"),
        LibraryTerm("lap_rho", lambda fd: fd.lap_rho(),
                    latex=r"\Delta\rho", equation="rho"),
    ]

def rho_terms_morse():
    return [
        LibraryTerm("div_rho_gradPhi", lambda fd: fd.div_rho_gradPhi(),
                    latex=r"\nabla\cdot(\rho\nabla\Phi)", equation="rho"),
    ]
\end{lstlisting}


\subsection{Default Library Builder}

\begin{lstlisting}[caption={multifield.py --- build\_default\_library}]
def build_default_library(*, morse=True, rich=False):
    """Build the recommended thesis-grade library."""
    rho_lib = rho_terms_core()
    if morse:
        rho_lib += rho_terms_morse()
    rho_lib += rho_terms_nonlinear_diff()
    if rich:
        rho_lib += rho_terms_coupling()

    px_lib = (px_terms_linear() + px_terms_pressure()
              + px_terms_viscosity())
    if morse:
        px_lib += px_terms_morse()
    if rich:
        px_lib += (px_terms_selfadvect()
                   + px_terms_hyperviscosity()
                   + px_terms_pressure_extended())

    py_lib = ...  # symmetric to px
    return {"rho": rho_lib, "px": px_lib, "py": py_lib}
\end{lstlisting}


% ============================================================================
\section{Multi-Field Weak System}
\label{sec:multifield_system}
% ============================================================================

For the 3-equation system, the weak form is built \emph{per equation}.
Unlike the scalar case (where the kernel encodes the operator via
integration by parts), each multi-field library term \emph{directly
provides} its full $(T,N_y,N_x)$ field.  The convolution is always with
the plain test function $\psi$ (identity kernel):

\begin{equation}\label{eq:multifield_weak}
  b_q^{(\alpha)} = (\psi_t * u_\alpha)(t_q, x_q, y_q),
  \qquad
  G_{q,m}^{(\alpha)} = (\psi * \Theta_m^{(\alpha)})(t_q, x_q, y_q),
\end{equation}
where $\alpha \in \{\rho, p_x, p_y\}$ and $\Theta_m^{(\alpha)}$ is the
$m$-th library term field for equation $\alpha$.

\begin{lstlisting}[caption={multifield.py --- Multi-field weak system}]
def build_weak_system_multifield(fd, psi_bundle, library,
                                  query_idx, target_field):
    """Build b=Gw for one equation of the multi-field PDE."""
    periodic = (False, True, True)
    qt, qx, qy = query_idx[:, 0], query_idx[:, 1], query_idx[:, 2]
    K, M = query_idx.shape[0], len(library)

    # LHS: b = conv(psi_t, target)
    psi_t = psi_bundle["psi_t"]
    b_full = fft_convolve3d_same(target_field, psi_t, periodic)
    b = b_full[qt, qx, qy]

    # RHS: each term convolved with psi (identity kernel)
    psi_kern = psi_bundle["psi"]
    G = np.empty((K, M))
    col_names = []
    for m, term in enumerate(library):
        field_m = term.evaluator(fd)   # (T, ny, nx) - already has derivatives
        conv_m = fft_convolve3d_same(field_m, psi_kern, periodic)
        G[:, m] = conv_m[qt, qx, qy]
        col_names.append(term.name)

    return b, G, col_names
\end{lstlisting}

\paragraph{Stacking across trajectories.}
Multiple training trajectories $\{(\rho^{(i)}, \bp^{(i)})\}$ are
stacked vertically:
\begin{equation}
  \mathbf{b} = \begin{pmatrix} \mathbf{b}^{(1)} \\ \vdots \\ \mathbf{b}^{(N_{\rm traj})} \end{pmatrix},
  \qquad
  \mathbf{G} = \begin{pmatrix} \mathbf{G}^{(1)} \\ \vdots \\ \mathbf{G}^{(N_{\rm traj})} \end{pmatrix}.
\end{equation}


% ============================================================================
\section{Regression: MSTLS}
\label{sec:regression}
% ============================================================================

\subsection{Preconditioning}

Before regression, each column of $\mathbf{G}$ is normalised to unit
$\ell_2$ norm:
\begin{equation}
  \tilde{G}_{:,m} = \frac{G_{:,m}}{\|G_{:,m}\|_2},
  \qquad
  s_m = \|G_{:,m}\|_2.
\end{equation}
After fitting $\tilde{\mathbf{w}}$, the physical coefficients are
recovered by $w_m = \tilde{w}_m / s_m$.

\begin{lstlisting}[caption={regression.py --- Column preconditioning}]
def precondition_columns(G):
    """Scale each column to unit L2 norm, return (G_scaled, scales)."""
    norms = np.linalg.norm(G, axis=0)
    norms[norms == 0] = 1.0
    return G / norms, norms
\end{lstlisting}


\subsection{Dominant-Balance Thresholding}

The key innovation of MSTLS over standard STRidge is the
\emph{dominant-balance mask}: a term is kept only if its normalised
contribution lies within a threshold band:
\begin{equation}\label{eq:dominant_balance}
  \hat{\lambda} \;\le\;
  \frac{|w_m|\,\|G_{:,m}\|}{\|\mathbf{b}\|}
  \;\le\; \frac{1}{\hat{\lambda}}.
\end{equation}
Terms whose contribution is too small (noise) or too large (dominant
but trivial) are pruned.

\begin{lstlisting}[caption={regression.py --- Dominant-balance mask}]
def _dominant_balance_mask(w, G, b, lam_hat):
    """Boolean mask: True = term passes dominant-balance filter."""
    b_norm = np.linalg.norm(b)
    if b_norm == 0:
        return np.ones(len(w), dtype=bool)

    col_norms = np.linalg.norm(G, axis=0)
    ratio = np.abs(w) * col_norms / b_norm

    return (ratio >= lam_hat) & (ratio <= 1.0 / lam_hat)
\end{lstlisting}


\subsection{MSTLS Algorithm}

The full MSTLS procedure sweeps over a grid of threshold values
$\hat{\lambda} \in \{\hat{\lambda}_1, \ldots, \hat{\lambda}_L\}$:

\begin{enumerate}[leftmargin=*]
\item For each $\hat{\lambda}$:
  \begin{enumerate}
    \item Solve $\tilde{\mathbf{w}} = \argmin \|\tilde{\mathbf{G}}\mathbf{w} - \mathbf{b}\|_2$
      (SVD least squares).
    \item Compute the dominant-balance mask.
    \item Restrict to active columns and re-solve.
    \item Iterate until convergence (active set stabilises or max iterations).
  \end{enumerate}
\item Compute the \emph{normalised loss} for each $\hat{\lambda}$:
  \begin{equation}
    L(\hat{\lambda}) = \frac{\|\mathbf{G}_{\rm active}\,\mathbf{w}_{\rm active} - \mathbf{b}\|_2}{\|\mathbf{b}\|_2}.
  \end{equation}
\item Select the $\hat{\lambda}$ minimising $L$, with ties broken by sparsity.
\end{enumerate}

\begin{lstlisting}[caption={regression.py --- MSTLS main loop}]
def mstls(G, b, lambdas, max_iter=25, tol=0.0):
    """Modified Sequential Threshold Least Squares.

    Returns best (w, active_mask, best_lambda, diagnostics).
    """
    K, M = G.shape
    b_norm = np.linalg.norm(b)

    best_w = None
    best_loss = np.inf
    best_lambda = lambdas[0]
    best_active = np.ones(M, dtype=bool)

    for lam in lambdas:
        active = np.ones(M, dtype=bool)
        for _ in range(max_iter):
            idx = np.where(active)[0]
            if len(idx) == 0:
                break
            w_sub = solve_ls(G[:, idx], b)  # SVD least-squares
            w_full = np.zeros(M)
            w_full[idx] = w_sub
            new_active = _dominant_balance_mask(w_full, G, b, lam)
            new_active &= (w_full != 0)
            if np.array_equal(new_active, active):
                break
            active = new_active

        # Compute normalised loss
        idx = np.where(active)[0]
        if len(idx) > 0:
            w_sub = solve_ls(G[:, idx], b)
            residual = G[:, idx] @ w_sub - b
            loss = np.linalg.norm(residual) / max(b_norm, 1e-15)
        else:
            loss = 1.0

        if loss < best_loss or (loss == best_loss
                and active.sum() < best_active.sum()):
            best_loss = loss
            best_lambda = lam
            best_active = active.copy()
            best_w = np.zeros(M)
            if len(idx) > 0:
                best_w[idx] = w_sub

    r2 = 1.0 - best_loss**2
    return best_w, best_active, best_lambda, {"normalised_loss": best_loss,
                                               "r2": r2}
\end{lstlisting}


% ============================================================================
\section{Fitting Pipeline}
\label{sec:fit}
% ============================================================================

The high-level fit (\texttt{fit.py}) combines preconditioning, MSTLS,
and coefficient unscaling:

\begin{lstlisting}[caption={fit.py --- wsindy\_fit\_regression}]
def wsindy_fit_regression(b, G, col_names, lambdas=None, max_iter=25,
                          tol=0.0):
    """Precondition -> MSTLS -> unscale -> WSINDyModel."""
    if lambdas is None:
        lambdas = np.logspace(-4, 1, 40)

    G_sc, scales = precondition_columns(G)
    w_sc, active, best_lam, diag = mstls(G_sc, b, lambdas,
                                          max_iter=max_iter, tol=tol)

    # Unscale: w_physical = w_scaled / scale
    w = np.zeros_like(w_sc)
    w[active] = w_sc[active] / scales[active]

    return WSINDyModel(
        col_names=list(col_names),
        w=w, active=active,
        best_lambda=best_lam,
        col_scale=scales,
        diagnostics=diag,
    )
\end{lstlisting}


% ============================================================================
\section{Model Selection}
\label{sec:select}
% ============================================================================

\subsection{Composite Score}

The test-function scale $\ell$ is a hyperparameter.  We sweep an
$\ell$-grid and score each trial by a composite criterion:
\begin{equation}\label{eq:composite_score}
  S(\ell) = L_{\rm norm}(\ell)
  + \alpha\,\frac{n_{\rm active}(\ell)}{M}
  + \beta\,\max\!\bigl(0,\;\log_{10}\kappa(\ell) - \log_{10}\kappa_0\bigr),
\end{equation}
where $L_{\rm norm}$ is the normalised loss, $n_{\rm active}/M$ is the
complexity ratio, $\kappa$ is the condition number of the active
sub-matrix, and $\alpha = 0.1$, $\beta = 0.01$, $\kappa_0 = 10^8$ are
defaults.

\paragraph{Pareto frontier.}
We also extract the Pareto frontier over $(L_{\rm norm}, n_{\rm active})$
as non-dominated solutions.


\subsection{Multi-Field Model Selection}

For the 3-equation system, each $\ell$ produces three models.
The score is:
\begin{equation}
  S_{\rm multi}(\ell) = \frac{1}{3}\bigl(R^2_\rho + R^2_{p_x} + R^2_{p_y}\bigr)
  - 0.05\,\frac{n_{\rm active}}{n_{\rm lib}}.
\end{equation}

\begin{lstlisting}[caption={multifield.py --- model\_selection\_multifield}]
def model_selection_multifield(field_list, library, ell_grid, ...):
    best_score = -np.inf
    for ell in ell_grid:
        result = discover_multifield(field_list, library, ell, ...)
        r2_rho = result.rho_model.diagnostics.get("r2", 0)
        r2_px  = result.px_model.diagnostics.get("r2", 0)
        r2_py  = result.py_model.diagnostics.get("r2", 0)
        n_act  = (result.rho_model.n_active
                  + result.px_model.n_active
                  + result.py_model.n_active)
        n_lib  = sum(len(library[eq]) for eq in library)
        avg_r2 = (r2_rho + r2_px + r2_py) / 3
        score = avg_r2 - 0.05 * n_act / max(n_lib, 1)

        if score > best_score:
            best_score = score
            best_result = result
    return best_result, best_ell
\end{lstlisting}


% ============================================================================
\section{WSINDy Model}
\label{sec:model}
% ============================================================================

\begin{lstlisting}[caption={model.py --- WSINDyModel dataclass}]
@dataclass
class WSINDyModel:
    col_names: List[str]        # Library term names
    w: np.ndarray               # Coefficient vector (M,)
    active: np.ndarray          # Boolean mask (M,)
    best_lambda: float          # Selected threshold
    col_scale: np.ndarray       # Column norms used for preconditioning
    diagnostics: Dict[str, Any] # r2, normalised_loss, etc.

    @property
    def n_active(self):
        return int(self.active.sum())

    @property
    def active_terms(self):
        return [n for n, a in zip(self.col_names, self.active) if a]

    def summary(self):
        lines = [f"WSINDyModel: {self.n_active} active terms"]
        for i, (name, w) in enumerate(zip(self.col_names, self.w)):
            if self.active[i]:
                lines.append(f"  {w:+12.4e}  {name}")
        return "\n".join(lines)
\end{lstlisting}


% ============================================================================
\section{Strong-Form RHS Evaluation}
\label{sec:rhs}
% ============================================================================

For time integration, we need to evaluate $u_t = \sum_m w_m\,\mathcal{D}_m[f_m(u)]$
in strong form (pointwise on the grid):

\begin{lstlisting}[caption={rhs.py --- wsindy\_rhs}]
def wsindy_rhs(u, grid, model, Lx=None, Ly=None):
    """Evaluate discovered RHS in strong form using spectral derivatives."""
    rhs = np.zeros_like(u)
    for i in range(len(model.w)):
        if not model.active[i]:
            continue

        op, feat = parse_term(model.col_names[i])
        f_u = eval_feature_pointwise(u, feat)   # u, u^2, u^3, ...

        # Apply differential operator
        if op == "I":
            rhs += model.w[i] * f_u
        elif op == "dx":
            rhs += model.w[i] * spectral_dx(f_u, Lx)
        elif op == "dy":
            rhs += model.w[i] * spectral_dy(f_u, Ly)
        elif op == "lap":
            rhs += model.w[i] * spectral_lap(f_u, Lx, Ly)
    return rhs
\end{lstlisting}

For the multi-field system, each equation's RHS is evaluated using the
\texttt{FieldData} container with spectral derivatives:

\begin{lstlisting}[caption={multifield.py --- multifield\_rhs}]
def multifield_rhs(rho, px, py, grid, Lx, Ly, rho_model, px_model,
                   py_model, rho_terms, px_terms, py_terms,
                   morse_params=None, ...):
    """Evaluate RHS of discovered 3-equation system."""
    # Build single-frame FieldData (T=1), spectral mode
    fd = FieldData(rho[np.newaxis], px[np.newaxis], py[np.newaxis],
                   grid, Lx, Ly, Phi=Phi, use_spectral=True)

    def _eval_rhs(model, terms, fd):
        rhs = np.zeros((fd.shape[1], fd.shape[2]))
        for i, t in enumerate(terms):
            if model.active[i]:
                rhs += model.w[i] * t.evaluator(fd)[0]
        return rhs

    return _eval_rhs(rho_model, rho_terms, fd), \
           _eval_rhs(px_model, px_terms, fd), \
           _eval_rhs(py_model, py_terms, fd)
\end{lstlisting}


% ============================================================================
\section{Time Integration}
\label{sec:integrators}
% ============================================================================

\subsection{Classical RK4}

For non-stiff systems:
\begin{align}
  k_1 &= \Delta t\, F(u^n), \\
  k_2 &= \Delta t\, F\bigl(u^n + \tfrac{1}{2}k_1\bigr), \\
  k_3 &= \Delta t\, F\bigl(u^n + \tfrac{1}{2}k_2\bigr), \\
  k_4 &= \Delta t\, F\bigl(u^n + k_3\bigr), \\
  u^{n+1} &= u^n + \tfrac{1}{6}\bigl(k_1 + 2k_2 + 2k_3 + k_4\bigr).
\end{align}


\subsection{ETDRK4: Exponential Time Differencing}

For stiff systems with a linear part $u_t = Lu + N(u)$:

\paragraph{Linear operator in Fourier space.}
We identify terms that are \emph{linear in} $u$ and map their spatial
operators to Fourier multipliers:
\begin{equation}\label{eq:L_hat}
  \hat{L}(k_x, k_y) = \sum_{m \in \text{linear}} w_m\,\sigma_m(k_x, k_y),
\end{equation}
where the symbol map is:
\begin{center}
\begin{tabular}{ll}
\toprule
$\mathcal{D}_m$ & $\sigma_m(k)$ \\
\midrule
$I$ & $1$ \\
$\partial_x$ & $ik_x$ \\
$\Delta$ & $-(k_x^2 + k_y^2)$ \\
$\Delta^2$ & $(k_x^2 + k_y^2)^2$ \\
\bottomrule
\end{tabular}
\end{center}

\begin{lstlisting}[caption={multifield.py --- extracting $\hat{L}$ for ETDRK4}]
def _extract_linear_spectral(model, terms, nx, ny, Lx, Ly):
    """Extract L_hat (ny, nx) from linear terms, or None."""
    K2 = _spectral_wavenumbers_2d(ny, nx, Lx, Ly)
    L_hat = np.zeros((ny, nx), dtype=np.complex128)
    has_linear = False
    for i, t in enumerate(terms):
        if not model.active[i]:
            continue
        c = model.w[i]
        if t.name.startswith("lap_") and "rho" not in t.name:
            L_hat += c * (-K2)
            has_linear = True
        elif t.name in ("px", "py"):
            L_hat += c   # linear decay/growth
            has_linear = True
    return L_hat if has_linear else None
\end{lstlisting}


\paragraph{ETDRK4 coefficients (Kassam \& Trefethen 2005).}
Given $z = \hat{L}\,\Delta t$, the $\varphi$-functions are computed via
contour integrals for numerical stability:
\begin{align}
  \varphi_1(z) &= \frac{e^z - 1}{z}, \\
  \varphi_2(z) &= \frac{e^z - 1 - z}{z^2}, \\
  \varphi_3(z) &= \frac{e^z - 1 - z - z^2/2}{z^3},
\end{align}
approximated by averaging over $M = 32$ points on a circle in the
complex plane:
\begin{equation}
  \varphi_k(z) \approx \frac{1}{M}\sum_{j=1}^M \varphi_k(z + r_j),
  \qquad
  r_j = e^{2\pi i(j-1/2)/M}.
\end{equation}

\begin{lstlisting}[caption={multifield.py --- ETDRK4 coefficient computation}]
def _etdrk4_coefficients(L_hat, dt):
    M = 32  # contour integration points
    z = L_hat * dt
    r = np.exp(2j * np.pi * (np.arange(1, M+1) - 0.5) / M)
    z_3d = z[..., np.newaxis]
    zr = z_3d + r          # (ny, nx, M)

    E  = np.exp(z)
    E2 = np.exp(z / 2)

    phi1   = np.mean((np.exp(zr) - 1) / zr, axis=-1)
    phi2   = np.mean((np.exp(zr) - 1 - zr) / zr**2, axis=-1)
    phi3   = np.mean((np.exp(zr) - 1 - zr - 0.5*zr**2) / zr**3, axis=-1)
    phi1_h = np.mean((np.exp(zr/2) - 1) / (zr/2), axis=-1)  # half-step

    a21 = dt * phi1_h
    b1 = dt * (4*phi3 - 3*phi2 + phi1)
    b2 = dt * 2 * (phi2 - 2*phi3)
    b3 = b2.copy()
    b4 = dt * (-phi2 + 4*phi3)

    return E, E2, a21, b1, b2, b3, b4
\end{lstlisting}


\paragraph{ETDRK4 time-stepping.}
Each step computes 4 stages in Fourier space:
\begin{align}
  \hat{a} &= e^{z/2}\,\hat{u}^n + a_{21}\,\hat{N}_1, \\
  \hat{b} &= e^{z/2}\,\hat{u}^n + a_{21}\,\hat{N}_2, \\
  \hat{c} &= e^{z/2}\,\hat{a} + a_{21}\,(2\hat{N}_3 - \hat{N}_1), \\
  \hat{u}^{n+1} &= e^z\,\hat{u}^n + b_1\hat{N}_1 + b_2(\hat{N}_2+\hat{N}_3) + b_4\hat{N}_4,
\end{align}
where $\hat{N}_j = \text{FFT}[N(u_j)]$ at the intermediate states.


\subsection{Physical Projections}

After each time step, two projections enforce physical consistency:
\begin{enumerate}
  \item \textbf{Non-negativity}: $\rho \leftarrow \max(\rho, 0)$
  \item \textbf{Mass conservation}: $\rho \leftarrow \rho \cdot M_0 / \!\int\!\rho\,dx$
\end{enumerate}

\begin{lstlisting}[caption={multifield.py --- physical projections}]
def _project(rho, px_cur, py_cur):
    if clip_negative_rho:
        rho = np.maximum(rho, 0.0)
    if mass_conserve and M0 > 0:
        M = float(np.sum(rho))
        if M > 0:
            rho = rho * (M0 / M)
    return rho, px_cur, py_cur
\end{lstlisting}


\subsection{Forecast Dispatch}

\begin{lstlisting}[caption={multifield.py --- forecast\_multifield (dispatch)}]
def forecast_multifield(rho0, px0, py0, result, grid, Lx, Ly,
                        n_steps, morse_params=None, ...
                        method="auto"):
    """
    method="auto": use ETDRK4 if any equation has Laplacian/
    bilaplacian terms, else RK4.
    """
    # Extract linear operators
    L_rho = _extract_linear_spectral(result.rho_model, ...)
    L_px  = _extract_linear_spectral(result.px_model, ...)
    L_py  = _extract_linear_spectral(result.py_model, ...)

    has_stiff = any(L is not None for L in [L_rho, L_px, L_py])
    use_etdrk4 = (method == "etdrk4") or (method == "auto" and has_stiff)

    if not use_etdrk4:
        # Classical RK4 loop
        for n in range(n_steps):
            k1r, k1x, k1y = rhs(rho, px, py)
            ...
            rho, px, py = _project(rho, px, py)
    else:
        # ETDRK4 loop in Fourier space
        coeff_rho = _etdrk4_coefficients(L_rho, dt)
        ...
        for n in range(n_steps):
            # 4 stages of ETDRK4 per field
            ...
\end{lstlisting}


% ============================================================================
\section{Bootstrap Uncertainty Quantification}
\label{sec:bootstrap}
% ============================================================================

\subsection{Scalar Bootstrap}

We assess coefficient uncertainty by resampling the query-point rows of
the weak system:
\begin{enumerate}
\item Draw $K$ rows with replacement from $(\mathbf{b}, \mathbf{G})$.
\item Refit MSTLS on the bootstrap sample.
\item Record $\mathbf{w}^{(b)}$ and which terms are active.
\item Repeat $B$ times (default $B=100$).
\end{enumerate}

\begin{equation}
  P_{\rm incl}(m) = \frac{1}{B}\sum_{b=1}^B \mathbf{1}[w_m^{(b)} \ne 0],
  \qquad
  \text{CI}_{95\%}(m) = \bigl[w_m^{(2.5\%)},\; w_m^{(97.5\%)}\bigr].
\end{equation}

\begin{lstlisting}[caption={uncertainty.py --- bootstrap\_wsindy}]
def bootstrap_wsindy(b, G, col_names, B=100, *, lambdas=None,
                     rng=None, replace=True, subsample_frac=1.0):
    rng = rng or np.random.default_rng()
    K, M = G.shape
    n_sample = max(1, int(K * subsample_frac))

    base_model = wsindy_fit_regression(b, G, col_names, lambdas=lambdas)
    coeff_samples = np.zeros((B, M))
    active_counts = np.zeros(M, dtype=np.intp)

    for rep in range(B):
        idx = rng.choice(K, size=n_sample, replace=replace)
        try:
            m = wsindy_fit_regression(b[idx], G[idx, :], col_names,
                                       lambdas=lambdas)
            coeff_samples[rep, :] = m.w
            active_counts += m.active.astype(np.intp)
        except Exception:
            pass

    return BootstrapResult(coeff_samples, active_counts, col_names,
                            base_model, B)
\end{lstlisting}


\subsection{Multi-Field Bootstrap}

For the 3-equation system, bootstrap operates at the \emph{trajectory
level}: resample entire trajectories (not individual rows) to capture
inter-trajectory variability.

\begin{lstlisting}[caption={multifield.py --- bootstrap\_multifield}]
def bootstrap_multifield(field_list, library, ell, ..., B=50, seed=42):
    rng = np.random.default_rng(seed)
    N_traj = len(field_list)
    for eq_name, target_fn in [("rho", ...), ("px", ...), ("py", ...)]:
        for rep in range(B):
            # Resample *trajectories* (not rows)
            idx = rng.choice(N_traj, size=N_traj, replace=True)
            fd_sub = [field_list[i] for i in idx]
            model, _, _, _ = fit_equation_multifield(fd_sub, ...)
            samples[rep] = model.w
            active_counts += model.active
\end{lstlisting}


% ============================================================================
\section{Full Discovery Pipeline}
\label{sec:pipeline}
% ============================================================================

Putting it all together, the \texttt{discover\_multifield} function:

\begin{enumerate}[leftmargin=*]
\item \textbf{Build fields:}
  $\texttt{build\_field\_data}(\rho, \text{traj}, \text{vel}, \ldots) \to \texttt{FieldData}$
  (one per trajectory, containing $\rho, \bp, \Phi$ and all cached derivatives).

\item \textbf{Build library:}
  $\texttt{build\_default\_library}(\text{morse}=\text{True}) \to
  \{\rho: [\text{terms}], p_x: [\text{terms}], p_y: [\text{terms}]\}$.

\item \textbf{Fit each equation:}
  For each $\alpha \in \{\rho, p_x, p_y\}$:
  \begin{enumerate}
    \item Construct test function $\psi$ at scale $\ell$.
    \item Stack weak systems across trajectories.
    \item Precondition $\to$ MSTLS $\to$ unscale.
    \item Record \texttt{WSINDyModel}.
  \end{enumerate}

\item \textbf{Model selection:}
  Sweep $\ell$-grid, score by average $R^2_{\rm weak}$ minus complexity penalty.

\item \textbf{Forecast:}
  Auto-detect stiffness $\to$ ETDRK4 or RK4 $\to$ time-step with
  projections (non-negativity, mass conservation).

\item \textbf{Uncertainty:}
  Bootstrap at trajectory level $\to$ coefficient distributions,
  inclusion probabilities, 95\% CIs.
\end{enumerate}

\begin{lstlisting}[caption={multifield.py --- discover\_multifield}]
def discover_multifield(field_list, library, ell, p=(2,2,2),
                        stride=(2,2,2), lambdas=None, verbose=True):
    if lambdas is None:
        lambdas = np.logspace(-4, 1, 40)

    results = {}
    for eq_name, target_fn in [("rho", lambda fd: fd.rho),
                                ("px", lambda fd: fd.px),
                                ("py", lambda fd: fd.py)]:
        model, b, G, cn = fit_equation_multifield(
            field_list, library[eq_name], target_fn,
            ell=ell, p=p, stride=stride, lambdas=lambdas)
        results[eq_name] = (model, b, G, cn)

    return MultiFieldResult(
        rho_model=results["rho"][0],
        px_model=results["px"][0],
        py_model=results["py"][0],
        rho_terms=library["rho"],
        px_terms=library["px"],
        py_terms=library["py"],
    )
\end{lstlisting}


% ============================================================================
\section{Standalone Spectral Operators}
\label{sec:operators}
% ============================================================================

The \texttt{operators.py} module provides spectral derivatives as
standalone functions (used before \texttt{FieldData} was introduced):

\begin{lstlisting}[caption={operators.py --- standalone spectral utilities}]
def fft_wavenumbers(n, d):
    """Wavenumbers: k_j = 2*pi*fftfreq(n, d)."""
    return 2 * np.pi * np.fft.fftfreq(n, d)

def grad_spectral(u, dx, dy):
    """Spectral gradient: (u_x, u_y)."""
    ny, nx = u.shape
    kx = fft_wavenumbers(nx, dx)
    ky = fft_wavenumbers(ny, dy)
    u_hat = np.fft.fft2(u)
    ux = np.real(np.fft.ifft2(1j * kx[np.newaxis, :] * u_hat))
    uy = np.real(np.fft.ifft2(1j * ky[:, np.newaxis] * u_hat))
    return ux, uy

def laplacian_spectral(u, dx, dy):
    """Spectral Laplacian: -(kx^2 + ky^2) * u_hat."""
    ny, nx = u.shape
    kx = fft_wavenumbers(nx, dx)
    ky = fft_wavenumbers(ny, dy)
    K2 = kx[np.newaxis, :]**2 + ky[:, np.newaxis]**2
    u_hat = np.fft.fft2(u)
    return np.real(np.fft.ifft2(-K2 * u_hat))
\end{lstlisting}


% ============================================================================
\section{Summary of Mathematical Objects}
\label{sec:math_summary}
% ============================================================================

\begin{center}
\renewcommand{\arraystretch}{1.3}
\begin{tabular}{lp{8.5cm}}
\toprule
\textbf{Symbol} & \textbf{Definition} \\
\midrule
$\Omega = [0,L_x)\times[0,L_y)$ & Periodic spatial domain \\
$\rho(x,y,t)$ & Density field (KDE from particles) \\
$\bp(x,y,t) = (p_x, p_y)$ & Polarisation / flux field \\
$W(r) = -C_r e^{-r/l_r} + C_a e^{-r/l_a}$ & Morse interaction kernel \\
$\Phi = W * \rho$ & Morse potential (FFT convolution) \\
$\psi = \varphi_t\,\varphi_x\,\varphi_y$ & Separable test function \\
$\varphi_i(s) = (1 - s^2/(\ell_i\Delta_i)^2)_+^{p_i}$ & 1-D polynomial bump \\
$\ell = (\ell_t, \ell_x, \ell_y)$ & Test-function half-widths \\
$p = (p_t, p_x, p_y)$ & Polynomial exponents \\
$\mathbf{b} = \mathbf{G}\mathbf{w}$ & Weak-form linear system \\
$b_q = (\psi_t * u)(q)$ & LHS at query point $q$ \\
$G_{q,m} = (\psi^{(m)} * f_m(u))(q)$ & RHS column $m$ at query $q$ \\
$\hat{\lambda}$ & Dominant-balance threshold \\
$L_{\rm norm} = \|G_{\rm act} w - b\|/\|b\|$ & Normalised residual \\
$\hat{L}(k) = \sum_m w_m\,\sigma_m(k)$ & Linear operator in Fourier space \\
$\varphi_k(z) = \frac{e^z - \sum_{j<k} z^j/j!}{z^k}$ & ETDRK4 $\varphi$-functions \\
$M_0 = \int\rho_0\,dx$ & Conserved total mass \\
\bottomrule
\end{tabular}
\end{center}


\end{document}
